\documentclass[a4paper, 14pt]{extarticle}

% Поля
%--------------------------------------
\usepackage{geometry}
\geometry{a4paper,tmargin=2cm,bmargin=2cm,lmargin=3cm,rmargin=1cm}
%--------------------------------------

\usepackage{float}
%Russian-specific packages
%--------------------------------------
\usepackage[T2A]{fontenc}
\usepackage[utf8]{inputenc} 
\usepackage[english, main=russian]{babel}
%--------------------------------------

\usepackage{textcomp}

% Красная строка
%--------------------------------------
\usepackage{indentfirst}               
%--------------------------------------             
\usepackage[most]{tcolorbox}
\tcbuselibrary{listings,breakable}
\newtcblisting{breakablelisting}{
  listing only,
  breakable,
  enhanced,
  frame hidden,
  colback=white,
  listing options={
    mathescape=false,
    texcl=false,
    inputencoding=utf8,
    literate={\$}{{\$}}1
  }
}


%Graphics
%--------------------------------------
\usepackage{graphicx}
\usepackage{wrapfig}
%--------------------------------------
\graphicspath{{./lab5/}}
% Полуторный интервал
%--------------------------------------
\linespread{1.3}                    
%--------------------------------------

%Выравнивание и переносы
%--------------------------------------
% Избавляемся от переполнений
\sloppy
% Запрещаем разрыв страницы после первой строки абзаца
\clubpenalty=10000
% Запрещаем разрыв страницы после последней строки абзаца
\widowpenalty=10000
%--------------------------------------

%Списки
\usepackage{enumitem}

%Подписи
\usepackage{caption} 

%Гиперссылки
\usepackage{hyperref}

\hypersetup {
	unicode=true
}

%Рисунки
%--------------------------------------
\DeclareCaptionLabelSeparator*{emdash}{~--- }
\captionsetup[figure]{labelsep=emdash,font=onehalfspacing,position=bottom}
%--------------------------------------

\usepackage{tempora}

%Листинги
%--------------------------------------
\usepackage{listings}
\lstset{
  basicstyle=\ttfamily\footnotesize, 
  %basicstyle=\footnotesize\AnkaCoder,        % the size of the fonts that are used for the code
  breakatwhitespace=false,         % sets if automatic breaks shoulbd only happen at whitespace
  breaklines=true,                 % sets automatic line breaking
  captionpos=t,                    % sets the caption-position to bottom
  inputencoding=utf8,
  frame=single,                    % adds a frame around the code
  keepspaces=true,                 % keeps spaces in text, useful for keeping indentation of code (possibly needs columns=flexible)
  keywordstyle=\bf,       % keyword style
  numbers=left,                    % where to put the line-numbers; possible values are (none, left, right)
  numbersep=5pt,                   % how far the line-numbers are from the code
  xleftmargin=25pt,
  xrightmargin=25pt,
  showspaces=false,                % show spaces everywhere adding particular underscores; it overrides 'showstringspaces'
  showstringspaces=false,          % underline spaces within strings only
  showtabs=false,                  % show tabs within strings adding particular underscores
  stepnumber=1,                    % the step between two line-numbers. If it's 1, each line will be numbered
  tabsize=2,                       % sets default tabsize to 8 spaces
  title=\lstname                   % show the filename of files included with \lstinputlisting; also try caption instead of title
}
%--------------------------------------

%%% Математические пакеты %%%
%--------------------------------------
\usepackage{amsthm,amsfonts,amsmath,amssymb,amscd}  % Математические дополнения от AMS
\usepackage{mathtools}                              % Добавляет окружение multlined
\usepackage[perpage]{footmisc}
%--------------------------------------

%--------------------------------------
%			НАЧАЛО ДОКУМЕНТА
%--------------------------------------

\begin{document}

%--------------------------------------
%			ТИТУЛЬНЫЙ ЛИСТ
%--------------------------------------
\begin{titlepage}
\thispagestyle{empty}
\newpage


%Шапка титульного листа
%--------------------------------------
\vspace*{-60pt}
\hspace{-65pt}
\begin{minipage}{0.3\textwidth}
\hspace*{-20pt}\centering
\includegraphics[width=\textwidth]{emblem}
\end{minipage}
\begin{minipage}{0.67\textwidth}\small \textbf{
\vspace*{-0.7ex}
\hspace*{-6pt}\centerline{Министерство науки и высшего образования Российской Федерации}
\vspace*{-0.7ex}
\centerline{Федеральное государственное автономное образовательное учреждение }
\vspace*{-0.7ex}
\centerline{высшего образования}
\vspace*{-0.7ex}
\centerline{<<Московский государственный технический университет}
\vspace*{-0.7ex}
\centerline{имени Н.Э. Баумана}
\vspace*{-0.7ex}
\centerline{(национальный исследовательский университет)>>}
\vspace*{-0.7ex}
\centerline{(МГТУ им. Н.Э. Баумана)}}
\end{minipage}
%--------------------------------------

%Полосы
%--------------------------------------
\vspace{-25pt}
\hspace{-35pt}\rule{\textwidth}{2.3pt}

\vspace*{-20.3pt}
\hspace{-35pt}\rule{\textwidth}{0.4pt}
%--------------------------------------

\vspace{1.5ex}
\hspace{-35pt} \noindent \small ФАКУЛЬТЕТ\hspace{80pt} <<Информатика и системы управления>>

\vspace*{-16pt}
\hspace{47pt}\rule{0.83\textwidth}{0.4pt}

\vspace{0.5ex}
\hspace{-35pt} \noindent \small КАФЕДРА\hspace{50pt} <<Теоретическая информатика и компьютерные технологии>>

\vspace*{-16pt}
\hspace{30pt}\rule{0.866\textwidth}{0.4pt}
  
\vspace{11em}

\begin{center}
\Large {\bf Лабораторная работа № 4.1} \\ 
\large {\bf по курсу <<Методы оптимизации>>} \\
\large <<Метод Хука и Дживса в 3D>> 
\end{center}\normalsize

\vspace{8em}


\begin{flushright}
  {Студент группы ИУ9-82Б Шемякин В.А. \hspace*{15pt}\\ 
  \vspace{2ex}
  Преподаватель Посевин Д. П.\hspace*{15pt}}
\end{flushright}

\bigskip

\vfill
 

\begin{center}
\textsl{Москва 2026}
\end{center}
\end{titlepage}
%--------------------------------------
%		КОНЕЦ ТИТУЛЬНОГО ЛИСТА
%--------------------------------------

\renewcommand{\ttdefault}{pcr}

\setlength{\tabcolsep}{3pt}
\newpage
\setcounter{page}{2}

\section{Задание}\label{Sect::task}
Реализовать два варианта метода Хука и Дживса (с однократным движением по направлению и с многократным движением по направлению).

Траекторию движения вдоль минимизирующей последовательности отобразить в трехмерном пространстве, а значение целевой функции цветом. Значение целевой функции цветом показывается так: в начальной точке минимизирующей последовательности цвет траектории светлый, а при приближении к минимуму цвет траектории становится более темным.
\section{Результаты}\label{Sect::res}

Исходный код программы представлен в листинге 1.

\begin{breakablelisting}
using Plots
using Colors

rosenbrock(x) = sum(100 * (x[i+1] - x[i]^2)^2 + (x[i] - 1)^2 for i in 1:length(x)-1)
schwefel(x) = 418.9829 * length(x) - sum(xi -> xi * sin(sqrt(abs(xi))), x)
rastrigin(x) = 10 * length(x) + sum(xi -> xi^2 - 10 * cos(2π * xi), x)

function explore(f, x, h)
    n = length(x)
    x_new = copy(x)
    probes = Tuple{Vector{Float64},Bool}[]
    for i in 1:n
        f_cur = f(x_new)
        p = copy(x_new); p[i] += h[i]
        if f(p) < f_cur
            push!(probes, (p, true))
            x_new[i] += h[i]
        else
            push!(probes, (p, false))
            p2 = copy(x_new); p2[i] -= h[i]
            if f(p2) < f_cur
                push!(probes, (p2, true))
                x_new[i] -= h[i]
            else
                push!(probes, (p2, false))
            end
        end
    end
    return x_new, probes
end

function search_single(f, x0, h0, eps; alpha=2.0, beta=0.5, max_iter=10000)
    n = length(x0)
    x_base = copy(x0)
    h = isa(h0, Number) ? fill(Float64(h0), n) : copy(h0)
    traj = [copy(x_base)]
    frames = []
    iter = 0

    while maximum(abs.(h)) > eps && iter < max_iter
        x_e, probes = explore(f, x_base, h)
        dest = f(x_e) < f(x_base) ? copy(x_e) : copy(x_base)
        push!(frames, (base=copy(x_base), probes=probes, tl=length(traj), dest=dest))

        if f(x_e) < f(x_base)
            x_pat = x_e + alpha * (x_e - x_base)
            x_base = x_e
            push!(traj, copy(x_base))

            x_pe, probes2 = explore(f, x_pat, h)
            dest2 = f(x_pe) < f(x_base) ? copy(x_pe) : copy(x_base)
            push!(frames, (base=copy(x_pat), probes=probes2, tl=length(traj), dest=dest2))

            if f(x_pe) < f(x_base)
                x_base = x_pe
                push!(traj, copy(x_base))
            end
        else
            h .*= beta
        end
        iter += 1
    end

    return x_base, f(x_base), traj, frames, iter
end

function search_multi(f, x0, h0, eps; alpha=2.0, beta=0.5, max_iter=10000)
    n = length(x0)
    x_base = copy(x0)
    h = isa(h0, Number) ? fill(Float64(h0), n) : copy(h0)
    traj = [copy(x_base)]
    iter = 0

    while maximum(abs.(h)) > eps && iter < max_iter
        x_e, _ = explore(f, x_base, h)

        if f(x_e) < f(x_base)
            d = x_e - x_base
            x_base = x_e
            push!(traj, copy(x_base))

            while iter < max_iter
                x_pat = x_base + alpha * d
                x_pe, _ = explore(f, x_pat, h)
                if f(x_pe) < f(x_base)
                    x_prev = x_base
                    x_base = x_pe
                    d = x_base - x_prev
                    push!(traj, copy(x_base))
                    iter += 1
                else
                    break
                end
            end
        else
            h .*= beta
        end
        iter += 1
    end

    return x_base, f(x_base), traj, iter
end

function seg_color_green(t)
    RGB(0.6 * (1 - t), 0.4 + 0.6 * (1 - t), 0.6 * (1 - t))
end
function seg_color_blue(t)
    RGB(0.6 * (1 - t), 0.6 * (1 - t), 0.5 + 0.5 * (1 - t))
end

function draw_3d_one(f, traj, title_str, filename; green=true)
    xs = [p[1] for p in traj]
    ys = [p[2] for p in traj]
    zs = [p[3] for p in traj]
    fv = [f(p) for p in traj]
    rng = max(maximum(fv) - minimum(fv), 1e-10)
    span_x = maximum(xs) - minimum(xs)
    span_y = maximum(ys) - minimum(ys)
    span_z = maximum(zs) - minimum(zs)
    pad = 0.2 * max(span_x, span_y, span_z, 0.5)
    xlim = (minimum(xs) - pad, maximum(xs) + pad)
    ylim = (minimum(ys) - pad, maximum(ys) + pad)
    zlim = (minimum(zs) - pad, maximum(zs) + pad)
    plt = plot(legend=false, xlabel="x₁", ylabel="x₂", zlabel="x₃", title=title_str, xlims=xlim, ylims=ylim, zlims=zlim)
    for i in 1:length(traj)-1
        t = (fv[i] - minimum(fv)) / rng
        col = green ? seg_color_green(t) : seg_color_blue(t)
        plot!(plt, xs[i:i+1], ys[i:i+1], zs[i:i+1], seriestype=:path3d, color=col, lw=2)
    end
    savefig(plt, filename)
    display(plt)
end

function make_gif_3d_one(f, traj, title_str, filename; green=true, n_frames=80)
    L = length(traj)
    idx = L > 1 ? round.(Int, range(1, L, length=n_frames)) : fill(1, n_frames)
    xs = [p[1] for p in traj]
    ys = [p[2] for p in traj]
    zs = [p[3] for p in traj]
    fv = [f(p) for p in traj]
    rng = max(maximum(fv) - minimum(fv), 1e-10)
    span_x = maximum(xs) - minimum(xs)
    span_y = maximum(ys) - minimum(ys)
    span_z = maximum(zs) - minimum(zs)
    pad = 0.2 * max(span_x, span_y, span_z, 0.5)
    xlim = (minimum(xs) - pad, maximum(xs) + pad)
    ylim = (minimum(ys) - pad, maximum(ys) + pad)
    zlim = (minimum(zs) - pad, maximum(zs) + pad)
    anim = @animate for k in 1:n_frames
        i = idx[k]
        plt = plot(legend=false, title=title_str, xlabel="x₁", ylabel="x₂", zlabel="x₃", xlims=xlim, ylims=ylim, zlims=zlim)
        if i >= 2
            for j in 1:i-1
                t = (fv[j] - minimum(fv)) / rng
                col = green ? seg_color_green(t) : seg_color_blue(t)
                plot!(plt, xs[j:j+1], ys[j:j+1], zs[j:j+1], seriestype=:path3d, color=col, lw=2)
            end
            ti = (fv[i] - minimum(fv)) / rng
            col = green ? seg_color_green(ti) : seg_color_blue(ti)
            scatter!(plt, [xs[i]], [ys[i]], [zs[i]], color=col, ms=5, marker=:circle)
        else
            col = green ? seg_color_green(0) : seg_color_blue(0)
            scatter!(plt, [xs[1]], [ys[1]], [zs[1]], color=col, ms=5, marker=:circle)
        end
    end
    gif(anim, filename, fps=12)
end

const show_plots = "--plot" in ARGS || "-p" in ARGS
const show_anim = "--anim" in ARGS || "-a" in ARGS
eps = 1e-6

funcs = [
    ("Розенброк", rosenbrock, [0.0, 0.0, 0.0], 0.4),
    ("Швефель", schwefel, [350.0, 400.0, 450.0], 0.5),
    ("Растригин", rastrigin, [2.5, -1.5, 1.0], 0.3)
]

results = []

for (name, f, x0, h0) in funcs
    x1, fv1, t1, fr1, it1 = search_single(f, x0, h0, eps)
    x2, fv2, t2, it2 = search_multi(f, x0, h0, eps)

    println("\n$name:")
    println("  Однократный поиск по направлению:")
    println("    x* = ", round.(x1, digits=6))
    println("    f(x*) = ", round(fv1, digits=10))
    println("    итераций: ", it1)
    println("    точек в релаксационной последовательности: ", length(t1))
    println("  Многократный поиск по направлению:")
    println("    x* = ", round.(x2, digits=6))
    println("    f(x*) = ", round(fv2, digits=10))
    println("    итераций: ", it2)
    println("    точек в релаксационной последовательности: ", length(t2))

    push!(results, (name=name, f=f, t1=t1, t2=t2, x0=x0))
end

println("\nСравнение точек в релаксационной последовательности:")
for r in results
    println("  $(r.name): однократный = $(length(r.t1)),  многократный = $(length(r.t2))")
end

if show_plots
    draw_3d_one(rosenbrock, results[1].t1, "Розенброк 3D: однократный поиск", "lab4_rosenbrock_single.png", green=true)
    draw_3d_one(rosenbrock, results[1].t2, "Розенброк 3D: многократный поиск", "lab4_rosenbrock_multi.png", green=false)
    draw_3d_one(schwefel, results[2].t1, "Швефель 3D: однократный поиск", "lab4_schwefel_single.png", green=true)
    draw_3d_one(schwefel, results[2].t2, "Швефель 3D: многократный поиск", "lab4_schwefel_multi.png", green=false)
    draw_3d_one(rastrigin, results[3].t1, "Растригин 3D: однократный поиск", "lab4_rastrigin_single.png", green=true)
    draw_3d_one(rastrigin, results[3].t2, "Растригин 3D: многократный поиск", "lab4_rastrigin_multi.png", green=false)
    if show_anim
        make_gif_3d_one(rosenbrock, results[1].t1, "Розенброк 3D: однократный поиск", "lab4_rosenbrock_single.gif", green=true)
        make_gif_3d_one(rosenbrock, results[1].t2, "Розенброк 3D: многократный поиск", "lab4_rosenbrock_multi.gif", green=false)
        make_gif_3d_one(schwefel, results[2].t1, "Швефель 3D: однократный поиск", "lab4_schwefel_single.gif", green=true)
        make_gif_3d_one(schwefel, results[2].t2, "Швефель 3D: многократный поиск", "lab4_schwefel_multi.gif", green=false)
        make_gif_3d_one(rastrigin, results[3].t1, "Растригин 3D: однократный поиск", "lab4_rastrigin_single.gif", green=true)
        make_gif_3d_one(rastrigin, results[3].t2, "Растригин 3D: многократный поиск", "lab4_rastrigin_multi.gif", green=false)
    end
end
\end{breakablelisting}

\newpage

Результат представлен на рисунках~\ref{fig:img1}--~\ref{fig:rastrigin}.

\begin{figure}[H]
  \centering
  \includegraphics[width=0.6\textwidth]{lab4_1/screen1.png}
  \caption{Результат}
  \label{fig:img1}
\end{figure}

\begin{figure}[H]
  \centering
  \includegraphics[width=0.6\textwidth]{lab4_1/screen2.png}
  \caption{Результат}
  \label{fig:schwefel}
\end{figure}

\begin{figure}[H]
  \centering
  \includegraphics[width=0.6\textwidth]{lab4_1/screen3.png}
  \caption{Результат}
  \label{fig:rastrigin}
\end{figure}

\end{document}
